\section{Descrição dos Dados}
Este trabalho utiliza três conjuntos de dados públicos: Energy Efficiency, Wine Quality (Red) e Iris. 

O \textbf{Energy Efficiency Dataset} contém 768 amostras de edifícios simulados, com 8 atributos construtivos (compactação, área, altura, orientação, etc.) e duas variáveis alvo: carga térmica de aquecimento (Y1) e de resfriamento (Y2). Todos os atributos são numéricos e não há valores ausentes.

O \textbf{Wine Quality Dataset (Red)} possui 1599 amostras de vinhos tintos portugueses, com 11 atributos físico-químicos (acidez, açúcar, pH, álcool, etc.) e uma variável alvo: nota de qualidade (0 a 10). Os dados são numéricos e desbalanceados entre as classes.

O \textbf{Iris Dataset} é composto por 150 amostras de flores, com 4 atributos morfológicos (comprimento/largura de sépala e pétala) e uma variável alvo (espécie). Não há valores ausentes e as classes são balanceadas.
